%:doctype: book
%:title: The Embroidermodder Reference Manual
%:author: The Embroidermodder Team
%:email: embroidermodder@gmail.com
%:description: Reference Manual for the technical details of all of the Embroidermodder Projects
%:url-repo: \url{https://github.com/embroidermodder/www.libembroidery.org}
%:toc: left
%:toclevels: 2
%:bibtex-file: references.bib
%:bibtex-style: apa
%:bibtex-order: alphabetical

%include:revision.adoc}

%// Only display the menu on the website.
%ifeval{["{backend}" == "html5"}
%include{menu.adoc}
%endif{}

\chapter{Introduction}

Reference Manual. For
Embroidermodder 2.0.0-alpha4, libembroidery 1.0.0-alpha, PET 1.0.0-alpha
and EmbroideryMobile 1.0.0-alpha.

Since the document is shipped automatically try to update the revnumber each
time you edit using `revision.sh`.

Test these:

\begin{verbatim}
sudo apt install latexml texlive-latex-base imagemagick info2man

# For our command line tools:
makeinfo embroider.texi -o embroider.info
info2man embroider.info > embroider.1
texi2pdf embroider.texi
# Or groff macro package for example ms.
# These may be housed in libembroidery since they're to be shipped as part of
# the embroider tarball.

# For online documentation:
pandoc embroidermodder_refman.tex -f latex -t html -s -o emb_refman.html --bibliography embroidermodder.bib
# Or latexml/latexmlpost
\end{verbatim}

\subsection{Command Language}

Printer Command Language (PCL), see \citet{packard1992pcl}.

HP-GL/2 Vector Graphics  index{HP-GL/2} described in \citet{packard1992pcl}.
Has commands like: "PU" Pen Up, "PR" Plot Relative, "EP" edge polygon.

So commands read like this:

\begin{verbatim}
PA40,10;
\end{verbatim}

command argument seperator(\texttt{,}) argument terminator(\texttt{;})

Constructing new commands from old ones in the command language is less
natural in the HP-GL/2 language, but a similar layer for us is
the tajima DST format \citep{4} for existing printers and CNC commands for
direct control... where'd we'd use G-Code \citep{7} and Linux CNC \citep{6}.

Could we `setpagedevice` to a printer in some cases and a similar CUPS service
for embroidery machines in others?

All systems are supported by ghostscript, when you account for homebrew \citet{9}:

\begin{verbatim}
brew update
brew upgrade
brew install ghostscript
brew cleanup
\end{verbatim}

Vector graphic logos don't require the SVG Qt library.

\subsection{Man Pages}

We maintain a traditional manpage for \texttt{embroider} using
the basic macros.

\subsection{Arduino}

\begin{verbatim}
apt-get install avr-libc gcc-avr uisp avrdude
\end{verbatim}

\section{Libembroidery}

(Under construction, please wait for v1.0 release.)

Libembroidery is a low-level library for reading, writing,
and altering digital embroidery files in C. It is part of the Embroidermodder Project
for open source machine embroidery.

Libembroidery is the underlying library that is used by Embroidermodder 2
and is developed by  The Embroidermodder Team \ref{the-embroidermodder-team}.
A full list of contributors to the project is maintained in
https://github.com/Embroidermodder/embroidermodder} the Embroidermodder 2 github
in the file `CREDITS.md`.
It handles over 45 different embroidery specific formats as well
as several non-embroidery specific vector formats.

It also includes a CLI called `embroider` that allows for better automation of
changes to embroidery files and will be more up-to date than
the Embroidermodder 2 GUI.

\subsection{Documentation}

Libembroidery is documented as part of this reference manual. If you need
libembroidery for any non-trivial usage or want to contribute to the library we
advise you read the appropriate design sections of the manual first. Copies of
this manual will be shipped with the packaged version of libembroidery, but to
build it we use the Doxyfile in
https://github.com/Embroidermodder/embroidermodder} the Embroidermodder git
repository.

For more basic usage, `embroider` should have some in-built help
starting with:

----
$ embroider --help
----

\subsection{License}

Libembroidery is distributed under the permissive zlib licence, see the LICENCE
file.

\section{Demos}

We're currently trying out some fill techniques which will be demonstrated here
and in the script `qa_test.sh`.

// center and caption
image{images/examples/logo.png[the logo}

Converts to:

// center and caption
image:images/examples/crossstitch_logo.png[crossstitch logo}

\subsection{Build}

libembroidery and EmbroiderModder 2 use CMake builds
so if you are building the project to use as a library we recommend
you run:

\begin{verbatim}
git clone https://github.com/Embroidermodder/libembroidery
cd libembroidery
cmake .
cmake --build .
cmake --install .
\end{verbatim}

This builds both the static and shared versions of the library as well
as the command line program `embroider`.

\citep{packard1992pcl}
\citep{linuxcncsrc}
\citep{linuxcnc}
\citep{adobe1990postscript}
\citep{postscript1999postscript}
\citep{eduTechDST}
\citep{cups}
\citep{millOperatorsManual}
\citep{oberg1914machinery}
\citep{dxf_reference}
\citep{embroidermodder_source_code}
\citep{libembroidery_source_code}
\citep{acatina}
\citep{kde_tajima}
\citep{wotsit_archive}
\citep{wotsit_siterip}
\citep{fineemb_dst}
\citep{edutechwiki_dst}

\section{Graphical User Interface for PC}
\ref{GUI}

\subsection{Overview}

(_UNDER MAJOR RESTRUCTURING, PLEASE WAIT FOR VERSION 2_)

\url{https://www.libembroidery.org}

Embroidermodder is a free machine embroidery application.
The newest version, Embroidermodder 2 can:

\begin{itemize}
\item edit and create embroidery designs
\item estimate the amount of thread and machine time needed to stitch a design
\item convert embroidery files to a variety of formats
\item upscale or downscale designs
\item run on Windows, Mac and Linux
\end{itemize}

\emph{Embroidermodder 2} is very much a work in progress since we're doing a ground
up rewrite to an interface in C using the GUI toolkit SDL2.
The reasoning for this is detailed in the issues tab.

For a more in-depth look at what we are developing read our
website (https://www.libembroidery.org}) which includes these docs as well
as the up-to date printer-friendly versions. These discuss recent changes,
plans and has user and developer guides for all the Embroidermodder projects.

To see what we're focussing on right now, see the Open Collective
News (https://opencollective.com/embroidermodder}).

// fixme
This current printer-friendly version
is here (https://www.libembroidery.org/EM2.0.0-alpha_refman_a4.pdf}).

\subsection{License}

The source code is under the terms of the zlib license: see `LICENSE.md`
in the source code directory.

Permission is granted to copy, distribute and/or modify this document
under the terms of the GNU Free Documentation License, Version 1.3
or any later version published by the Free Software Foundation;
with no Invariant Sections, no Front-Cover Texts, and no Back-Cover Texts.

A copy of the license is included in Section~\ref{GNU-free-documentation-license}.

\section{The Embroidermodder Project and Team}

The _Embroidermodder 2_ project is a collection of small software
utilities for manipulating, converting and creating embroidery files in all
major embroidery machine formats. The program _Embroidermodder 2_ itself
is a larger graphical user interface (GUI) which is at the heart of the project.

The tools and associated documents are:

\begin{itemize}
\item The website \url{https://www.libembroidery.org}
\item This reference manual covering the development of all these projects with
the current version available here:
\url{https://www.libembroidery.org/embroidermodder_2.0_manual.pdf}
\item The GUI _Embroidermodder 2_ covered in Chapter~\ref{GUI}.
\item The core library of low-level functions: `libembroidery`, covered in
Chapter~\ref{libembroidery}
\item The CLI `embroider`, which is part of libembroidery
\item Mobile embroidery format viewers and tools convered in Chapter~\ref{Mobile}.
\item Specs for an open source hardware embroidery machine extension called
the Portable Embroidery Tool (PET) which is also part of libembroidery. See
Chapter~\ref{PET}.
\end{itemize}

The website, this manual and some scripts to construct the distribution are
maintained in \citep{embroidermodder_website}.

They all tools to make the standard
user experience of working with an embroidery machine better without expensive
software which is locked to specific manufacturers and formats. But ultimately
we hope that the core _Embroidermodder 2_ is a practical, ever-present tool in
larger workshops, small cottage industry workshops and personal hobbyist's
bedrooms.

Embroidermodder 2 is licensed under the zlib license and we aim to keep all of our tools open
source and free of charge. If you would like to support the project check out our  Open
Collective (https://opencollective.com/embroidermodder}) group.
If you would like to help,
please join us on GitHub. This document is written as developer training as well helping new
users (see the last sections) so this is the place to learn how to start changing the code.

The Embroidermodder Team is the collection of people who've submitted
patches, artwork and documentation to our three projects.
The team was established by Jonathan Greig and Josh Varga.
The full list is actively maintained below.

\subsection{Core Development Team}

Embroidermodder 2:

\begin{itemize}
\item Jonathan Greig (\url{https://github.com/redteam316})
\item Josh Varga (\url{https://github.com/JoshVarga})
\item Robin Swift (\url{https://github.com/robin-swift})
\end{itemize}

Embroidermodder 1:

\begin{itemize}
\item Josh Varga (https://github.com/JoshVarga)
\item Mark Pontius (http://sourceforge.net/u/mpontius/profile)
\end{itemize}

\section{Credits for Embroidermodder 2, libembroidery and all other related code}

If you have contributed and wish to be added to this list, alter the  README on Embroidermodder
github page (https://github.com/Embroidermodder/Embroidermodder) and we'll copy it to the
libembroidery source code since that is credited to "The Embroidermodder Team".

\section{Embroidermodder 1}

The Embroidermodder Team is also inspired by the original Embroidermodder that
was built by Mark Pontius and the same Josh Varga on SourceForge which
unfortunately appears to have died from linkrot. We may create a distribution
on here to be the official "legacy" Embroidermodder code but likely in a
seperate repository because it's GNU GPL v3 and this code is written to be
zlib (that is, permissive licensed) all the way down.

One reason why this is useful is that the rewrite by Jonathan Greig, John Varga
and Robin Swift for Embroidermodder 2 should have no regressions: no features
present in v1 should be missing in v2.

\section{Features}

Embroidermodder 2 has many advanced features that enable you to create awesome designs quicker, tweak existing designs to perfection, and can be fully customized to fit your workflow.

A summary of these features:

\begin{itemize}
\item Cross Platform
\item Realistic rendering
\item Various grid types and auto-adjusting rulers
\item Many measurement tools
\item Add text to any design
\item Supports many formats
\item Batch Conversion
\item Scripting API
\end{itemize}

\subsection{Cross Platform}

If you use multiple operating systems, it's important to choose software that works on all of them.

Embroidermodder 2 runs on Windows, Linux and Mac OS X. Let's not forget the  Raspberry
Pi (https://www.raspberrypi.org}).

// [#running-different-platforms}
.Running on different platforms.
image:features-platforms-1.png[Running on different platforms}

\subsection{Realistic Rendering}

It is important to be able to visualize what a design will look like when stitched and our
pseudo ``3D'' realistic rendering helps achieve this
(see Figure~\ref{real-render}).

// [#fig-real-render}
.Real render examples.
image:features-realrender-1.png}
image:features-realrender-2.png}
image:features-realrender-3.png}

\subsection{Various grid types and auto-adjusting rulers}

Making use of the automatically adjusting ruler in conjunction with the grid will ensure your
design is properly sized and fits within your embroidery hoop area.

Use rectangular, circular or isometric grids to construct your masterpiece!

Multiple grids and rulers in action Figure ref fig grid-ruler.

// [#fig-grid-ruler}
.Grid and ruler examples.
image{features-grid-ruler-1.png}

\subsection{Many measurement tools}

Taking measurements is a critical part of creating great designs. Whether you are designing
mission critical embroidered space suits for NASA or some other far out design for your next
meet-up, you will have precise measurement tools at your command to make it happen. You can
locate individual points or find distances between any 2 points anywhere in the design!

Take quick and accurate measurements:

image{images/features-measure-1.png}

\subsection{Add text to any design}

Need to make company apparel for all of your employees with individual names on them? No sweat.
Just simply add text to your existing design or create one from scratch, quickly and easily.
Didn't get it the right size or made a typo? No problem. Just select the text and update it
with the property editor.

Add text and adjust its properties quickly:

image{images/features-text-1.png}

\subsection{Supports many formats}

Embroidery machines all accept different formats. There are so many formats
available that it can sometimes be confusing whether a design will work with your machine.

Embroidermodder 2 supports a wide variety of embroidery formats as well as several vector
formats, such as SVG and DXF. This allows you to worry less about which designs you can use.

\subsection{Batch Conversion}

Need to send a client several different formats? Just use libembroidery-convert, our command
line utility which supports batch file conversion.

There are a multitude of formats to choose from:

image{images/features-formats-1.png[features formats}

\subsection{Scripting API}

If you've got programming skills and there is a feature that isn't currently available that you
absolutely cannot live without, you have the capability to create your own custom commands for
Embroidermodder 2. We provide an QtScript API which exposes various application functionality
so that it is possible to extend the application without requiring a new release. If you have
created a command that you think is worth including in the next release, just  contact
us (contact.html) and we will review it for functionality, bugs, and finally inclusion.

An Embroidermodder 2 command excerpt:

image{images/features-scripting-1.png[scripting screenshot}

\subsection{Build and Install}

Assuming you already have the SDL2 libraries you can proceed to using the fast build, which
assumes you want to build and test locally.

The fast build should be:

\begin{verbatim}
bash build.sh
\end{verbatim}

or, on Windows:

\begin{verbatim}
.\build.bat
\end{verbatim}

Then run using the `run.bat` or `run.sh` scripts in the build/ directory.

Otherwise, follow the instructions below.

If you plan to install the dev version to your system (we recommend you wait
for the official installers and beta release first) then use the CMake build
instead.

\subsection{Install on Desktop}

We recommend that if you want to install the development version you use the CMake build. Like
this:

----
git submodule --init --update

mkdir build
cd build
cmake ..
cmake --build .
sudo cmake --install .
----

These lines are written into the file:

----
./build_install.sh
----

On Windows use the next section.

\section{History}

Embroidermodder 1 was started by Mark Pontius in 2004 while staying up all night with his son in his first couple months. When Mark returned
to his day job, he lacked the time to continue the project. Mark made the decision to focus on his family and work, and in
2005, Mark gave full control of the project to Josh Varga so that Embroidermodder could continue its growth.

Embroidermodder 2 was conceived in mid 2011 when Jonathan Greig and Josh Varga discussed the possibility of making a cross-platform version.
It is currently in active development and will run on GNU/Linux, Mac OS X, Microsoft Windows and Raspberry Pi.

All Embroidermodder downloads (downloads.html) are hosted on SourceForge.

The source code for Embroidermodder 1 (http://sourceforge.net/p/embroidermodder/code/HEAD/tree/embroidermodder1}) has always been hosted on Sourceforge.

The source code for Embroidermodder 2 (https://github.com/Embroidermodder/Embroidermodder}) was moved to GitHub on July 18, 2013.

The website for Embroidermodder (https://github.com/Embroidermodder/www.libembroidery.org}) was moved to GitHub on September 9, 2013.

\chapter{Contact us}

For general questions email:  embroidermodder at gmail.com (mailto:embroidermodder@gmail.com)

To request a new feature  open an issue on the main Embroidermodder GitHub repository (https://github.com/Embroidermodder/Embroidermodder/issues). We'll move it to the correct repository.

\chapter{Downloads}

\section{Alpha Build}

This is a highly experimental build: we recommend users wait for the beta release when the basic features
are functional.

Visit our  GitHub Releases page (https://github.com/Embroidermodder/Embroidermodder/releases)
for the current build. Unfortunately, earlier builds went down with the Sourceforge page we hosted them on.

\chapter{GUI}

Embroidermodder 2 is very much a work in progress since we're doing a ground up rewrite to an
interface in Python using the GUI toolkit Tk. The reasoning for this is detailed in the issues
tab.

For a more in-depth look at what we are developing read the developer notes (link to dev notes
section). This discusses recent changes in a less formal way than a changelog (since this
software is in development) and covers what we are about to try.

\section{Documentation}

The documentation is in the form of the website (included in the `docs/` directory) and the
printed docs in this file.

\subsection{Development}

If you wish to develop with us you can chat via the contact email on the
website (https://www.libembroidery.org}) or in the issues tab on the
github page (https://github.com/Embroidermodder/Embroidermodder/issues}).
People have been polite and friendly in these conversations and I (Robin) have
really enjoyed them. If we do have any arguments please note we have a Code of
Conduct (`CODE_OF_CONDUCT.md`) so there is a consistent policy to enforce when
dealing with these arguments.

The first thing you should try is building from source using the build advice(link to build)
above. Then read some of the development notes (link to dev notes.md) to get the general
layout of the source code and what we are currently planning.

\subsection{Testing}

To find unfixed errors run the tests by launching from the command line with:

\begin{lstlisting}
$ embroidermodder --test
\end{lstlisting}

then dig through the output. It's currently not worth reporting the errors, since
there are so many but if you can fix anything reported here you can submit a PR.

\section{Code Optimisations and Simplifications}

\subsection{Geometry}

The geometry is stored, processed and altered via libembroidery. See the Python specific part
of the documentation for libembroidery for this. What the code in Embroidermodder does is make
the GUI widgets to change and view this information graphically.

For example if we create a circle with radius 10mm and center at `(20mm, 30mm)` then fill it
with stitches the commands would be

\begin{lstlisting}
from libembroidery import Pattern, Circle, Vector, satin
circle = Circle(Vector(20, 30), 10)
pattern = Pattern()
pattern.add_circle(circle, fill=satin)
pattern.to_stitches()
\end{lstlisting}

but the user would do this through a series of GUI actions:

\begin{enumerate}
\item Create new file
\item Click add circle
\item Use the Settings dialog to alter the radius and center
\item Use the fill tool on circle
\item Select satin from the drop down menu
\end{enumerate}

So EM2 does the job of bridging that gap.

\subsection{Postscript Support}

In order to safely support user contributed/shared data that can define, for example, double to
double functions we need a consistent processor for these descriptions.

Embroidermodder backends to the postscript interpreter included in libembroidery to accomplish
this.

For example the string: `5 2 t mul add` is equivalent to
the expression $2*t + 5$.

The benefit of not allowing this to simply be a Python expression is that it is safe against
malicious use, or accidental misuse. The program can identify whether the output is of the
appropriate form and give finitely many calculations before declaring the function to have run
too long (stopping equations that hang).

To see examples of this see the `assets/shapes/*.ps` files.

\subsection{SVG Icons}

To make the images easier to alter and restyle we could switch to svg icons. There's some code
in the git history to help with this.

\subsection{The Actions System}

In order to simplify the development of a GUI that is flexible and easy to understand to new
developers we have a custom action system that all user actions will go via an `actuator` that
takes a string argument. By using a string argument the undo history is just an array of
strings.

The C \texttt{action\_hash\_data} struct will contain: the icon used, the
labels for the menus and tooltips and the function pointer for that action. There will be an
accompanying argument for this function call, currently being drafted as
`action_call`. So when the user makes a function call it should contain
information like the mouse position, whether special key is pressed etc.

\subsection{Accessibility}

Software can be more or less friendly to people with dyslexia, partial sightedness, reduced
mobility and those who don't speak English. Embroidermodder 2 has, in its design, the following
features to help:

\begin{itemize}
\item icons for everything to reduce the amount of reading required
\item the system font is configurable: if you have a dyslexia-friendly font you can load it
\item the interface rescales to help with partial-sightedness
\item the system language is configurable, unfortunately the docs will only be in English but we can try to supply lots of images of the interface to make it easier to understand as a second language
\item buttons are remappable: XBox controllers are known for being good for people with reduced mobility so remapping the buttons to whatever setup you have should help
\end{itemize}

Note that most of these features will be released with version 2.1, which is planned for around
early 2023.

\subsection{Sample Files}

Various sample embroidery design files can be found in the `embroidermodder2/samples` folder.

\subsection{Shortcuts}

A shortcut can be made up of zero or more modifier keys and at least one non-modifier key
pressed at once.

To make this list quickly assessable, we can produce a list of hashes which are simply the
flags ORed together.

The shortcuts are stored in the csv file `shortcuts.csv` as a 5-column table
with the first 4 columns describing the key combination. This is loaded into
the shortcuts `TABLE`. Each tick the program checks the input state for this
combination by first translating the key names into indices for the key state,
then checking for whether all of them are set to true.

\subsection{Removed Elements}

So I've had a few pieces of web infrastructure fail me recently and I think
it's worth noting. An issue that affects us is an issue that can effect people
who use our software.

\subsection{Qt and dependencies}

Downloading and installing Qt has been a pain for some users (46Gb on possibly
slow connections).

I'm switching to FreeGLUT 3 (which is a whole other conversation) which means
we can ship it with the source code package meaning only a basic build
environment is necessary to build it.

\subsection{Social Platform}

Github is giving me a server offline (500) error and is still giving a bad ping.

So... all the issues and project boards etc. being on Github is all well and
good assuming that we have our own copies. But we don't if Github goes down or
some other major player takes over the space and we have to move (again, since
this started on SourceForge).

This file is a backup for that which is why I'm repeating myself between them.

\subsection{OpenGL}

OpenGL rendering within the application. This will allow for Realistic Visualization - Bump
Mapping/OpenGL/Gradients?

This should backend to a C renderer or something.

\subsection{Configuration Data Ideas}

Embroidermodder should boot from the command line regardless of whether it is or is not
installed (this helps with testing and running on machines without root). Therefore, it can
create an initiation file but it won't rely on its existence to boot:
`~/.embroidermodder/config.json`.

* Switch colors to be stored as 6 digit hexcodes with a `#`.
* We've got close to a hand implemented ini read/write setup in `settings.py`.

\subsection{Distribution}
\index{distribution}

When we release the new pip wheel we should also package:

\begin{itemize}
\item `.tar.gz` and `.zip` source archive.
\item Debian package
\item RPM package
\end{itemize}

Only do this once per minor version number.

\begin{itemize}
\item todo Screenshot a working draft to demonstrate.
\end{itemize}

\subsection{Perennial Jobs}

* Check for memory leaks
* Clear compiler warnings on `-Wall -ansi -pedantic` for C.
* Write new tests for new code.
* Get Embroidermodder onto the current version of libembroidery.
* PEP7 compliance.
* Better documentation with more photos/screencaps.

\subsection{Full Test Suite}
\index{testing}

(This needs a hook from Embroidermodder to embroider's full test suite.)

The flag `--full-test-suite` runs all the tests that have been written.
Since this results in a lot of output the details are both to stdout
and to a text file called `test_matrix.txt`.

Patches that strictly improve the results in the `test_matrix.txt` over
the current version will likely be accepted and it'll be a good place
to go digging for contributions. (Note: strictly improve means that
the testing result for each test is as good a result, if not better.
Sacrificing one critera for another would require some design work
before we would consider it.)

\subsection{Symbols
index[symbols}

Symbols use the SVG path syntax.

In theory, we could combine the icons and symbols systems, since they could be
rendered once and stored as icons in Qt. (Or as textures in FreeGLUT.)

Also we want to render the patterns themselves using SVG syntax, so it would
save on repeated work overall.

\section{Features

\subsection{Bindings
index[bindings}

Bindings for libembroidery are maintained for the languages we use internally
in the project, for other languages we consider that the responsibility of
other teams using the library.

So libembroidery is going to be supported on:

* `C` (by default)
* `C++` (also by default)
* `Java` (for the \index{Android] application MobileViewer)
* `Swift` (for the \index{iOS] application iMobileViewer)

For `C\#` index[C-sharp] we recommend directly calling the function directly
using the DllImport feature:

----
/* Calling readCsv() via C# as a native function. */
[DllImport("libembroidery.so", EntryPoint="readCsv")}
----

see this StackOverflow discussion for help: https://stackoverflow.com/questions/11425202/is-it-possible-to-call-a-c-function-from-c-net}.

For Python you can do the same using ctypes: https://www.geeksforgeeks.org/how-to-call-a-c-function-in-python/}.

\subsection{Other Supported Thread Brands
index[supported threads}

The thread lists that aren't preprogrammed into formats but are indexed in
the data file for the purpose of conversion or fitting to images/graphics.

* Arc Polyester
* Arc Rayon
* Coats and Clark Rayon
* Exquisite Polyester
* Fufu Polyester
* Fufu Rayon
* Hemingworth Polyester
* Isacord Polyester
* Isafil Rayon
* Marathon Polyester
* Marathon Rayon
* Madeira Polyester
* Madeira Rayon
* Metro Polyester
* Pantone
* Robison Anton Polyester
* Robison Anton Rayon
* Sigma Polyester
* Sulky Rayon
* ThreadArt Rayon
* ThreadArt Polyester
* ThreaDelight Polyester
* Z102 Isacord Polyester

\section{GUI Design
index[GUI}

Embroidermodder 2 was written in C++/Qt5 and it was far too complex. We had issues with people
not able to build from source because the Qt5 libraries were so ungainly. So I decided to do a
rewrite in C/SDL2 (originally FreeGLUT, but that was a mistake) with data stored as YAML. This
means linking 4-7 libraries depending on your system which are all well supported and widely available.

This is going well, although it's slow progress as I'm trying to keep track of the design while
also doing a ground up rewrite. I don't want to throw away good ideas. Since I also write code
for libembroidery my time is divided.
Overview of the UI rewrite

(Problems to be solved in brackets.)

It's not much to look at because I'm trying to avoid using an external
widgets system, which in turn means writing things like toolbars and menubars
over. If you want to get the design the actuator is the heart of it.

Without Qt5 we need a way of assigning signals with actions, so this is what
I've got: the user interacts with a UI element, this sends an integer to the
actuator that does the thing using the current state of the mainwindow struct
of which we expect there to be exactly one instance. The action is taken out
by a jump table that calls the right function (most of which are missing in
action and not connected up properly). It also logs the number, along with
key parts of the main struct in the undo history (an unsolved problem because
we need to decide how much data to copy over per action). This means undo,
redo and repeat actions can refer to this data.

\section{To Do}

These should be sorted into relevant code sections.

\begin{itemize}
\item todo sort todo list.
\item Alpha: High priority
\begin{itemize}
\item Statistics from 1.0, needs histogram
\item Saving DST/PES/JEF (varga)
\item Saving CSV/SVG (rt) + CSV read/write UNKNOWN interpreted as COLOR bug
\end{itemize}
\item Alpha: medium priority
\begin{itemize}
\item Notify user of data loss if not saving to an object format.
\item Import Raster Image
\item SNAP/ORTHO/POLAR
\item Layer Manager + LayerSwitcher DockWidget
\item Reading DXF
\end{itemize}
\item Alpha: low priority
** Writing DXF
** Up and Down keys cycle thru commands in the command prompt
** Amount of Thread, Machine Time Estimation (also allow customizable times for setup, color changes, manually
trimming jump threads, etc...that way a realistic total time can be estimated)
** Otto Theme Icons - whatsthis icon doesn't scale well, needs redone
** embroidermodder2.ico 16 x 16 looks horrible
* Alpha: lower priority
** CAD Command: Arc (rt)
* beta
** Custom Filter Bug - doesn't save changes in some cases
** Cannot open file with `#` in name when opening multiple files (works fine when opening the single file)
** Closing Settings Dialog with the X in the window saves settings rather than discards them
** Advanced Printing
** Filling Algorithms (varga)
** Otto Theme Icons - beta (rt) - Units, Render, Selectors
* Finish before 2.0 release
** QDoc Comments
** Review KDE4 Thumbnailer
** Documentation for libembroidery and formats
** HTML Help files
** Update language translations
** CAD Command review: line
** CAD Command review: circle
** CAD Command review: rectangle
** CAD Command review: polygon
** CAD Command review: polyline
** CAD Command review: point
** CAD Command review: ellipse
** CAD Command review: arc
** CAD Command review: distance
** CAD Command review: locatepoint
** CAD Command review: move
** CAD Command review: rgb
** CAD Command review: rotate
** CAD Command review: scale
** CAD Command review: singlelinetext
** CAD Command review: star
** Clean up all compiler warning messages, right now theres plenty :P
* 2.0
** tar.gz archive
** zip archive
** Debian Package (rt)
** NSIS Installer (rt)
** Mac Bundle?
** press release
* 2.x/Ideas
** libembroidery.mk for MXE project (refer to qt submodule packages for qmake based building. Also refer to plibc.mk for example of how write an update macro for github.)
** libembroidery safeguard for all writers - check if the last stitch is an END stitch. If not, add an end stitch in the writer and modify the header data if necessary.
** Cut/Copy - Allow Post-selection
** CAD Command: Array
** CAD Command: Offset
** CAD Command: Extend
** CAD Command: Trim
** CAD Command: BreakAtPoint
** CAD Command: Break2Points
** CAD Command: Fillet
** CAD Command: Chamfer
** CAD Command: Split
** CAD Command: Area
** CAD Command: Time
** CAD Command: PickAdd
** CAD Command: Product
** CAD Command: Program
** CAD Command: ZoomFactor
** CAD Command: GripHot
** CAD Command: GripColor and GripCool
** CAD Command: GripSize
** CAD Command: Highlight
** CAD Command: Units
** CAD Command: Grid
** CAD Command: Find
** CAD Command: Divide
** CAD Command: ZoomWindow (Move out of view.cpp)
** Command: Web (Generates Spiderweb patterns)
** Command: Guilloche (Generates Guilloche patterns)
** Command: Celtic Knots
** Command: Knotted Wreath
** Lego Mindstorms NXT/EV3 ports and/or commands.
** native function that flashes the command prompt to get users attention when using the prompt is required for a command.
** libembroidery-composer like app that combines multiple files into one.
** Settings Dialog, it would be nice to have it notify you when switching tabs that a setting has been changed. Adding an Apply button is what would make sense for this to happen.
** Keyboard Zooming/Panning
** G-Code format?
** 3D Raised Embroidery
** Gradient Filling Algorithms
** Stitching Simulation
** RPM packages?
** Reports?
** Record and Playback Commands
** Settings option for reversing zoom scrolling direction
** Qt GUI for libembroidery-convert
** EPS format? Look at using Ghostscript as an optional add-on to libembroidery...
** optional compile option for including LGPL/GPL libs etc... with warning to user about license requirements.
** Realistic Visualization - Bump Mapping/OpenGL/Gradients?
** Stippling Fill
** User Designed Custom Fill
** Honeycomb Fill
** Hilbert Curve Fill
** Sierpinski Triangle fill
** Circle Grid Fill
** Spiral Fill
** Offset Fill
** Brick Fill
** Trim jumps over a certain length.
** FAQ about setting high number of jumps for more controlled trimming.
** Minimum stitch length option. (Many machines also have this option too)
** Add 'Design Details' functionality to libembroidery-convert
** Add 'Batch convert many to one format' functionality to libembroidery-convert
** EmbroideryFLOSS - Color picker that displays catalog numbers and names.
* beta
** Realistic Visualization - Bump Mapping/OpenGL/Gradients?
** Get undo history widget back (BUG).
** Mac Bundle, .tar.gz and .zip source archive.
** NSIS installer for Windows, Debian package, RPM package
** GUI frontend for embroider features that aren't
supported by embroidermodder: flag selector from a table
** Update all formats without color to check for edr or rgb files.
** Setting for reverse scrolling direction (for zoom, vertical pan)
** Keyboard zooming, panning
** New embroidermodder2.ico 16x16 logo that looks good at that scale.
** Saving dst, pes, jef.
** Settings dialog: notify when the user is switching tabs
that the setting has been changed, adding apply button is what would
make sense for this to happen.
** Update language translations.
** Replace KDE4 thumbnailer.
** Import raster image.
** Statistics from 1.0, needs histogram.
** SNAP/ORTHO/POLAR.
** Cut/copy allow post-selection.
** Layout into config.
** Notify user of data loss if not saving to an object format.
** Add which formats to work with to preferences.
** Cannot open file with `#` in the name when opening multiple
files but works with opening a single file.
** Closing settings dialog with the X in the window saves
settings rather than discarding them.
** Otto theme icons: units, render, selectors, what's
this icon doesn't scale.
** Layer manager and Layer switcher dock widget.
** Test that all formats read data in correct scale
(format details should match other programs).
** Custom filter bug -- doesn't save changes in some cases.
** Tools to find common problems in the source code and suggest fixes
to the developers. For example, a translation miss: that is, for any language
other than English a missing entry in the translation table should supply a
clear warning to developers.
** Converting Qt C++ version to native GUI C throughout.
** OpenGL Rendering: `Real` rendering to see what the embroidery
looks like, Icons and toolbars, Menu bar.
** Libembroidery interfacing: get all classes to use the proper
libembroidery types within them. So `Ellipse` has `EmbEllipse` as public
data within it.
** Move calculations of rotation and scaling into `EmbVector` calls.
** GUI frontend for embroider features that aren't supported by
embroidermodder: flag selector from a table
** Update all formats without color to check for edr or rgb files.
** Setting for reverse scrolling direction (for zoom, vertical pan)
** Keyboard zooming, panning
** Better integrated help: I don't think the help should backend to
a html file somewhere on the user's system. A better system would be a custom
widget within the program that's searchable.
** New embroidermodder2.ico 16x16 logo that looks good at that scale.
** Settings dialog: notify when the user is switching tabs that the
setting has been changed, adding apply button is what would make sense for
this to happen.
\end{itemize}

\section{Contributing}

\subsection{Version Control}

Being an open source project, developers can grab the latest code at any time
and attempt to build it themselves. We try our best to ensure that it will build smoothly
at any time, although occasionally we do break the build. In these instances,
please provide a patch, pull request which fixes the issue or open an issue and
notify us of the problem, as we may not be aware of it and we can build fine.

Try to group commits based on what they are related to: features/bugs/comments/graphics/commands/etc...

See the coding style  here (coding-style).

\subsection{Get the Development Build going}

When we switch to releases we recommend using them, unless you're reporting a bug in which case you can check the development build for whether it has been patched. If this applies to you, the current development build is https://github.com/Embroidermodder/Embroidermodder/releases/tag/alpha3[here].

\subsection{To Do}

\begin{itemize}
\item Beta
\begin{itemize}
\item Libembroidery 1.0.
\item Better integrated help: I don't think the help should backend to a html file somewhere on the user's system. A better system would be a custom widget within the program that's searchable.
\item EmbroideryFLOSS - Color picker that displays catalog numbers and names.
\item Custom filter bug -- doesn't save changes in some cases.
\item Advanced printing.
\item Stitching simulation.
\end{itemize}
\item 2.x/ideas
\begin{itemize}
\item User designed custom fill.
\end{itemize}
\end{itemize}

These are key bits of reasoning behind why the GUI is built the way it is.

\section{Translation of the user interface}

In a given table the left column is the default symbol and the right string is the translation.
If the translate function fails to find a translation it returns the default symbol.

So in US English it is an empty table, but in UK English
only the dialectical differences are present.

Ideally, we should support at least the 6 languages spoken at the UN. Quoting \url{https://www.un.org}

\begin{quote}
\emph{There are six official languages of the UN. These are Arabic, Chinese, English, French, Russian and Spanish.}
\end{quote}

We're adding Hindi, on the grounds that it is one of the most commonly spoken languages and at
least one of the Indian languages should be present.

Written Chinese is generally supported as two different symbol sets and we follow that
convension.

English is supported as two dialects to ensure that the development team is aware of what those
differences are. The code base is written by a mixture of US and UK native English speakers
meaning that only the variable names are consistently one dialect: US English. As for
documentation: it is whatever dialect the writer prefers (but they should maintain consistency
within a text block like this one).

Finally, we have "default", which is the dominant language
of the internals of the software. Practically, this is
just US English, but in terms of programming history this
is the "C locale".

\section{Old action system notes}

Action: the basic system to encode all user input.

This typedef gives structure to the data associated with each action
which, in the code, is referred to by the action id (an int from
the define table above).

\section{DESCRIPTION OF STRUCT CONTENTS}

\subsection{label}

The action label is always in US English, lowercase,
seperated with hyphens.

For example: \texttt{new-file}.

\section{Flags}

The bit based flags all collected into a 32-bit integer.

label{tab:flags-for-actions}
[cols="1,6",title="Flags of EM actions"}
\begin{tabular}{l l l}
|bit(s)
|description

|0
|User (0) or system (1) permissions.

|1-3
|The mode of input.

|4-8
|The object classes that this action can be applied to.

|9-10
|What menu (if any) should it be present in.

|11-12
|What
\end{tabular}\section{Description}

The string placed in the tooltip describing the action.

\section{Original Prompt System}

NOTE: `main()` is run every time the command is started.
Use it to reset variables so they are ready to go.

NOTE: `click()` is run only for left clicks.
Middle clicks are used for panning.
Right clicks bring up the context menu.

NOTE: `move()` is optional. It is run only after
`enableMoveRapidFire()` is called. It
will be called every time the mouse moves until
`disableMoveRapidFire()` is called.

NOTE: `prompt()` is run when Enter is pressed.
`appendPromptHistory` is automatically called before `prompt()`
is called so calling it is only needed for erroneous input.
Any text in the command prompt is sent as an uppercase string.

\section{CAD command review}

[cols="1,1,1,6"}
\begin{tabular}{l l l}
| **ID**
| **name**
| **arguments**
| **description**

| 0
| \index{\texttt{newfile}}
| none
| Create a new EmbPattern with a new tab in the GUI.

|1
|\index{\texttt{openfile}}
|filename string
|Open an EmbPattern with the supplied filename `fname`.

|2
|\index{\texttt{savefile}}
|filename string
|Save the current loaded EmbPattern to the supplied filname `fname`.

|3
|\index{\texttt{scale}}
|selected objects, 1 float
|Scale all selected objects by the number supplied, without selection scales the entire design.

|4
|\index{\texttt{circle}}
|mouse co-ords
|Adds a circle to the design based on the supplied numbers, converts to stitches on save for stitch only formats.

|5
|\index{\texttt{offset}}
|mouse co-ords
|Shifts the selected objects by the amount given by the mouse co-ordinates.

|6
|\index{\texttt{extend}}
|
|

|7
|\index{\texttt{trim}}
|
|

|8
|\index{\texttt{break_at_point}}
|
|

|9
|\index{\texttt{break_2_points}}
|
|

|10
|\index{\texttt{fillet}}
|
|

|11
|\index{\texttt{star}}
|
|

|12
|\index{\texttt{singlelinetext}}
|
|

|13
|\index{\texttt{chamfer}}
|
|

|14
|\index{\texttt{split}}
|
|

|15
|\index{\texttt{area}}
|
|

|16
|\index{\texttt{time}}
|
|

|17
|\index{\texttt{pickadd}}
|
|

|16
|\index{\texttt{zoomfactor}}
|
|

|17
|\index{\texttt{product}}
|
|

|18
|\index{\texttt{program}}
|
|

|19
|\index{\texttt{zoomwindow}}
|
|

|20
|\index{\texttt{divide}}
|
|

|21
|\index{\texttt{find}}
|
|

| 22
| \index{\texttt{record}}
|
|

| 23
| \index{\texttt{playback}}
|
|

| 24
| \index{\texttt{rotate}}
|
|

| 25
| \index{\texttt{rgb}}
|
|

| 26
| \index{\texttt{move}}
|
|

| 27
| \index{\texttt{grid}}
|
|

| 28
| \index{\texttt{griphot}}
|
|

| 29
| \index{\texttt{gripcolor}}
|
|

| 30
| \index{\texttt{gripcool}}
|
|

| 31
| \index{\texttt{gripsize}}
|
|

| 32
| \index{\texttt{highlight}}
|
|

| 33
| \index{\texttt{units}}
|
|

| 34
| \index{\texttt{locatepoint}}
|
|

| 35
| \index{\texttt{distance}}
|
|

| 36
| \index{\texttt{arc}}
|
|

| 37
| \index{\texttt{ellipse}}
|
|

|38
| \index{\texttt{array}}
|
|

|39
| \index{\texttt{point}}
|
|

|40
| \index{\texttt{polyline}}
|
|

|41
| \index{\texttt{polygon}}
|
|

| 42
| \index{\texttt{rectangle}}
|
|

| 43
| \index{\texttt{line}}
|
|

| 44
| \index{	exttt{arc-rt}}
|
|

| 45
| \index{\texttt{dolphin}}
|
|

| 46
| \index{\texttt{heart}}
|
|

\end{tabular}\section{Actions}

\subsection{ARC}
index[arc}

\subsection{CIRCLE}
index[circle}

\subsection{OPEN}
index[open}

\section{Changelog}

\section{Ideas}

Stuff that is now supposed to be generated by Doxygen:

* todo: Bibliography style to plainnat.
* todo: US letter paper version of printed docs.

\chapter{Formats}

\section{Overview}

\subsubsection{Read/Write Support Levels}

The table of read/write format support levels uses the status levels described here:

.Read/Write Support Levels
[cols="1,8"}
\begin{tabular}{l l l}
| *Status Label*
| *Description*

| `rw-none`
| Either the format produces no output, reporting an error. Or it produces a Tajima dst file as an alternative.

| `rw-poor`
| A file somewhat similar to our examples is produced. We don't know how well it runs on machines in practice as we don't have any user reports or personal tests.

| `rw-basic`
| Simple files in this format run well on machines that use this format.

| `rw-standard`
| Files with non-standard features work on machines and we have good documentation on the format.

| `rw-reliable`
| All known features don't cause crashes. Almost all work as expected.

| `rw-complete`
| All known features of the format work on machines that use this format. Translations from and to this format preserve all features present in both.
\end{tabular}These can be split into `r-basic w-none`, for example, if they don't match.

So all formats can, in principle, have good read and good write support, because it's defined in relation to files that we have described the formats for.

\subsubsection{Test Support Levels}

[cols="1,8"}
\begin{tabular}{l l l}
| *Status Label*
| *Description*

| `test-none`
| No tests have been written to test the specifics of the format.

| `test-basic`
| Stitch Lists and/or colors have read/write tests.

| `test-thorough`
| All features of that format has at least one test.

| `test-fuzz`
| Can test the format for uses of features that we haven't thought of by feeding in nonsense that is designed to push possibly dangerous weaknesses to reveal themselves.

| `test-complete`
| Both thorough and fuzz testing is covered.
\end{tabular}So all formats can, in principle, have complete testing support, because it's defined in relation to files that we have described the formats for.

\subsubsection{Documentation Support Levels}

% cols="1,8"
\begin{tabular}
\textbf{Status Label} &
\textbf{Description}
\\
| `doc-none` &
| We haven't researched this beyond finding example files.
\\
| `doc-basic` &
| We have a rough sketch of the size and contents of the header if there is one. We know the basic stitch encoding (if there is one), but not necessarily all stitch features.
\\
| `doc-standard` &
| We know some good sources and/or have tested all the features that appear to exist. They mostly work the way we have described.
\\
| `doc-good` &
| All features that were described somewhere have been covered here or we have thoroughly tested our ideas against other softwares and hardwares and they work as expected.
\\
| `doc-complete` &
| There is a known official description and our description covers all the same features.
\end{tabular}

Not all formats can have complete documentation because it's based on what
information is publically available. So the total score is reported in the table
below based on what level we think is available.

\subsubsection{Overall Support}

Since the overall support level is the combination of these
4 factors, but rather than summing up their values it's an
issue of the minimum support of the 4.

% cols="1,8"
\begin{tabular}
\textbf{Status Label} &
\textbf{Description}
\\
`read-only` &
If write support is none and read support is not none.
\\
`write-only` &
If read support is none and write support is not none.
\\
`unstable` &
If both read and write support are not none but testing or documentation is none.
\\
`basic` &
If all ratings are better than none.
\\
`reliable` &
If all ratings are better than basic.
\\
`complete` &
If all ratings could not reasonably be better (for example any improvements
rely on information that we may never have access to). This is the only status
that can be revoked, since if the format changes or new documentation is
released it is no longer complete.
\\
`experimental` &
For all other scenarios.
\end{tabular}

\section{Table of Format Support Levels}

Overview of documentation support by format.

[cols="1,1,1"}
\begin{tabular}{l l l}
| *Format*
| *Ratings*
| *Score*

| \index{Toyota] Embroidery Format (index[`100}}`.100`)
| rw-basic doc-none test-none
| unstable

| \index{Toyota] Embroidery Format (index[`10o}}`.10o`)
| rw-basic doc-none test-none
| unstable

| \index{Bernina] Embroidery Format (index[`art}}`.art`)
| rw-none doc-none test-none
| experimental

| \index{Bitmap Cache] Embroidery Format (`.bmc`)
| r-basic w-none doc-none test-none|unstable

| \index{Bits and Volts] Embroidery Format (`.bro`)
| rw-none doc-none test-none
| experimental

| Melco Embroidery Format (`.cnd`)
| rw-none doc-none test-none
| experimental

| Embroidery Thread Color Format (`.col`)
| rw-basic doc-none test-none
| `experimental`

| Singer Embroidery Format (`.csd`)
| rw-none doc-none test-none
| experimental

| Comma Separated Values (`.csv`)
| rw-none doc-none test-none
| experimental

| Barudan Embroidery Format (`.dat`)
| rw-none doc-none test-none
| experimental

| Melco Embroidery Format (.dem)
| rw-none doc-none test-none
| experimental

| Barudan Embroidery Format (.dsb)
| rw-none doc-none test-none
| experimental

| Tajima Embroidery Format (.dst)
| rw-none doc-none test-none
| experimental

| ZSK USA Embroidery Format (.dsz)
| rw-none doc-none test-none
| experimental

| Drawing Exchange Format (.dxf)
| rw-none doc-none test-none
| experimental

| Embird Embroidery Format (.edr)
| rw-none doc-none test-none
| experimental

| Elna Embroidery Format (.emd)
| rw-none doc-none test-none
| experimental

| Melco Embroidery Format (.exp)
| rw-none doc-none test-none
| experimental

| Eltac Embroidery Format (.exy)
| rw-none doc-none test-none
| experimental

| Sierra Expanded Embroidery Format (.eys)
| rw-none doc-none test-none
| experimental

| Fortron Embroidery Format (.fxy)
| rw-none doc-none test-none
| experimental

| Smoothie G-Code Embroidery Format (.gc)
| rw-none doc-none test-none
| experimental

| Great Notions Embroidery Format (.gnc)
| rw-none doc-none test-none
| experimental

| Gold Thread Embroidery Format (.gt)
| rw-none doc-none test-none
| experimental

| Husqvarna Viking Embroidery Format (.hus)
| rw-none doc-none test-none
| experimental

| Inbro Embroidery Format (.inb)
| rw-none doc-none test-none
| experimental

| Embroidery Color Format (.inf)
| rw-none doc-none test-none
| experimental

| Janome Embroidery Format (.jef)
| rw-none doc-none test-none
| experimental

| Pfaff Embroidery Format (.ksm)
| rw-none doc-none test-none
| experimental

| Pfaff Embroidery Format (.max)
| rw-none doc-none test-none
| experimental

| Mitsubishi Embroidery Format (.mit)
| rw-none doc-none test-none
| experimental

| Ameco Embroidery Format (.new)
| rw-none doc-none test-none
| experimental

| Melco Embroidery Format (`.ofm`)
| rw-none doc-none test-none
| experimental

| Pfaff Embroidery Format (.pcd)
| rw-none doc-none test-none
| experimental

| Pfaff Embroidery Format (`.pcm`)
| rw-none doc-none test-none
| experimental

| Pfaff Embroidery Format (`.pcq`)
| rw-none doc-none test-none
| experimental

| Pfaff Embroidery Format (`.pcs`)
| rw-none doc-none test-none
| experimental

| Brother Embroidery Format (`.pec`)
| rw-none doc-none test-none
| experimental

| Brother Embroidery Format (.pel)
| rw-none doc-none test-none
| experimental

| Brother Embroidery Format (.pem)
| rw-none doc-none test-none
| experimental

| Brother Embroidery Format (.pes)
| rw-none doc-none test-none
| experimental

| Brother Embroidery Format (.phb)
| rw-none doc-none test-none
| experimental

| Brother Embroidery Format (.phc)
| rw-none doc-none test-none
| experimental

| AutoCAD Embroidery Format (.plt)
| rw-none doc-none test-none
| experimental

| RGB Embroidery Format (.rgb)
| rw-none doc-none test-none
| experimental

| Janome Embroidery Format (.sew)
| rw-none doc-none test-none
| experimental

| Husqvarna Viking Embroidery Format (.shv)
| rw-none doc-none test-none
| experimental

| Sunstar Embroidery Format (.sst)
| rw-none doc-none test-none
| experimental

| Data Stitch Embroidery Format (.stx)
| rw-none doc-none test-none
| experimental

| Scalable Vector Graphics (.svg)
| rw-none doc-none test-none
| experimental

| Pfaff Embroidery Format (.t01)
| rw-none doc-none test-none
| experimental

| Pfaff Embroidery Format (.t09)
| rw-none doc-none test-none
| experimental

| Happy Embroidery Format (.tap)
| rw-none doc-none test-none
| experimental

| ThredWorks Embroidery Format (`.thr`)
| rw-none doc-none test-none
| experimental

| Text File (`.txt`)
| rw-none doc-none test-none
| experimental

| Barudan Embroidery Format (`.u00`)
| rw-none doc-none test-none
| experimental

| Barudan Embroidery Format (index[`u01}}`.u01`)
| rw-none doc-none test-none
| experimental

| Pfaff Embroidery Format (`.vip`)
| rw-none doc-none test-none
| experimental

| Pfaff Embroidery Format (`.vp3`)
| rw-none doc-none test-none
| experimental

| Singer Embroidery Format (`.xxx`)
| rw-none doc-none test-none
| experimental

| ZSK USA Embroidery Format (`.zsk`)
| rw-none doc-none test-none
| experimental

| *FORMAT*
| *READ, WRITE*
| *NOTES*

| \index{\texttt{10o}}
| YES, NO
| read (need to fix external color loading) (maybe find out what ctrl and code flags of `0x10`, `0x08`, `0x04`, and `0x02` mean)

| \index{\texttt{100}}
| NO, NO
| none (4 byte codes) `61 00 10 09` (type, type2, x, y ?) x and y (signed char)

| \index{\texttt{art}}
| NO, NO
| none

| \index{\texttt{bro}}
| YES NO
| read (complete)(maybe figure out detail of header)

| \index{\texttt{cnd}}
| NO, NO
| none

| \index{\texttt{col}}
| NO, NO
| (color file no design) read(final) write(final)

| \index{\texttt{csd}}
| YES NO
| read (complete)

| \index{\texttt{dat}}
| NO, NO
| read ()

| \index{\texttt{dem}}
| NO, NO
| none (looks like just encrypted cnd)

| \index{\texttt{dsb}}
| YES NO
| read (unknown how well) (stitch data looks same as 10o)

| \index{\texttt{dst}}
| YES NO
| read (complete) / write(unknown)

| \index{\texttt{dsz}}
| YES NO
| read (unknown)

| \index{\texttt{dxf}}
| NO, NO
| read (Port to C. needs refactored)

| \index{\texttt{edr}}
| NO, NO
| read (C version is broken) / write (complete)

| \index{\texttt{emd}}
| NO, NO
| read (unknown)

| \index{\texttt{exp}}
| YES NO
| read (unknown) / write(unknown)

| \index{\texttt{exy}}
| YES NO
| read (need to fix external color loading)

| \index{\texttt{fxy}}
| YES NO
| read (need to fix external color loading)

| \index{\texttt{gnc}}
| NO, NO
| none

| \index{\texttt{gt}}
| NO, NO
| read (need to fix external color loading)

| \index{\texttt{hus}}
| YES NO
| read (unknown) / write (C version is broken)

| \index{\texttt{inb}}
| YES NO
| read (buggy?)

| \index{\texttt{jef}}
| YES NO
| write (need to fix the offsets when it is moving to another spot)

| \index{\texttt{ksm}}
| YES NO
| read (unknown) / write (unknown)

| \index{\texttt{pcd}}
| NO, NO
|

| \index{\texttt{pcm}}
| NO, NO
|

| \index{\texttt{pcq}}
| NO, NO
| read (Port to C)

| \index{\texttt{pcs}}
| BUGGY, NO
| read (buggy / colors are not correct / after reading, writing any other format is messed up)

| \index{\texttt{pec}}
| NO, NO
| read / write (without embedded images, sometimes overlooks some stitches leaving a gap)

| \index{\texttt{pel}}
| NO, NO
| none

| \index{\texttt{pem}}
| NO, NO
| none

| \index{\texttt{pes}}
| YES, NO
|

| \index{\texttt{phb}}
| NO, NO
|

| \index{\texttt{phc}}
| NO, NO
|

| \index{\texttt{rgb}}
| NO, NO
|

| \index{\texttt{sew}}
| YES, NO
|

| \index{\texttt{shv}}
| NO, NO
| read (C version is broken)

| \index{\texttt{sst}}
| NO, NO
| none

| \index{\texttt{svg}}
| NO, YES
|

| \index{\texttt{tap}}
| YES, NO
| read (unknown)

| \index{\texttt{u01}}
| NO, NO
|

| \index{\texttt{vip}}
| YES, NO
|

| \index{\texttt{vp3}}
| YES, NO
|

| \index{\texttt{xxx}}
| YES, NO
|

| \index{\texttt{zsk}}
| NO, NO
| read (complete)
\end{tabular}* TODO Josh, Review this section and move any info still valid or
needing work into TODO comments in the actual
libembroidery code. Many items in this list are out of date and do not reflect the current status of
libembroidery. When finished, delete this file.
** Test that all formats read data in correct scale (format details should match other programs)
** Add which formats to work with to preferences.
** Check for memory leaks
** Update all formats without color to check for edr or rgb files
** Fix issues with DST (VERY important that DST work well)
* todoSupport for Singer FHE, CHE (Compucon) formats?

\chapter{Geometry and Algorithms

\section{To Do

\subsubsection{Arduino

* Fix emb-outline files
* Fix thread-color files
* Logging of Last Stitch Location to External USB Storage(commonly available and easily replaced) ...wait until TRE is available to avoid rework
* inotool.org - seems like the logical solution for Nightly/CI builds
* Smoothieboard experiments

\subsubsection{Testing

* looping test that reads 10 times while running valgrind. See `embPattern_loadExternalColorFile()` Arduino leak note for more info.

\subsubsection{Development

If you wish to develop with us you can chat via the contact email
on the  website https://libembroidery.org} or in the issues tab on the
 github page https://github.com/Embroidermodder/Embroidermodder/issues}.
People have been polite and friendly in these conversations and I (Robin)
have really enjoyed them.
If we do have any arguments please note we have a
 Code of Conduct  CODE\_OF\_CONDUCT.md so there is a consistent policy to
enforce when dealing with these arguments.

The first thing you should try is building from source using the  build advice (build)
above. Then read some of the  manual
https://libembroidery.org/embroidermodder_2.0_manual.pdf to get the general
layout of the source code and what we are currently planning.

\subsubsection{Testing

To find unfixed errors run the tests by launching from the command line with:

----
$ embroidermodder --test
----

then dig through the output. It's currently not worth reporting the errors, since
there are so many but if you can fix anything reported here you can submit a PR.

\section{Contributing

\subsubsection{Funding

The easiest way to help is to fund development (see the Donate button above),
since we can't afford to spend a lot of time developing and only have limited
kit to test out libembroidery on.

\subsubsection{Programming and Engineering

Should you want to get into the code itself:

* Low level C developers are be needed for the base library \texttt{libembroidery}.
* Low level assembly programmers are needed for translating some of \texttt{libembroidery} to \texttt{EmbroiderBot}.
* Hardware Engineers to help design our own kitbashed embroidery machine \texttt{EmbroiderBot}, one of the original project aims in 2013.
* Scheme developers and C/SDL developers to help build the GUI.
* Scheme developers to help add designs for generating of custom stitch-filled emblems like the heart or dolphi. Note that this happens in Embroidermodder not libembroidery (which assumes that you already have a function available).

\subsubsection{Writing

We also need people familiar with the software and the general
machine embroidery ecosystem to contribute to the
documentation (https://github.com/Embroidermodder/www.libembroidery.org).

We need researchers to find references for the documentation: colour tables,
machine specifications etc. The history is murky and often very poorly maintained
so if you know anything from working in the industry that you can share: it'd be
appreciated!

\section{Embroidermodder Project Coding Standards}

A basic set of guidelines to use when submitting code.

Code structure is mre important than style, so first we advise you read "Design" and experimenting before getting into the specifics of code style.

\subsubsection{Where Code Goes

Anything that deals with the specifics of embroidery file formats, threads,
rendering to images, embroidery machinery or command line interfaces should go 
in `libembroidery` not here.

\subsubsection{Where Non-compiled Files Go

TODO: Like most user interfaces Embroidermodder is mostly data, so here we will have a list describing where each CSV goes.}

\subsubsection{Ways in which we break style on purpose

Most style guides advise you to keep functions short. We make a few pointed
exceptions to this where the overall health and functionality of the source code should benefit.

The `actuator` function will always be a mess and it should be: we're keeping
the total source lines of code down by encoding all user action into a descrete
sequence of strings that are all below \texttt{\_STRING\_LENGTH} in length. See
the section on the actuator (TODO) describing why any other solution we could
think  here would mean more more code without a payoff in speed of execution or
clarity.

\section{Version Control}

Being an open source project, developers can grab the latest code at any time and attempt to build it themselves. We try our best to ensure that it will build smoothly at any time, although occasionally we do break the build. In these instances, please provide a patch, pull request which fixes the issue or open an issue and notify us of the problem, as we may not be aware of it and we can build fine.

Try to group commits based on what they are related to: features/bugs/comments/graphics/commands/etc...

\section{Donations

Creating software that interfaces with hardware is costly. A summary of some of the costs involved:

* Developer time for 2 core developers
* Computer equipment and parts
* Embroidery machinery
* Various electronics for kitbashing Embroiderbot
* Consumable materials (thread, fabric, stabilizer, etc...)

If you have found our software useful, please consider funding further development by donating to the project on Open Collective (https://opencollective.com/embroidermodder}).

\section{Embroidermodder Project Coding Standards

Rather than maintain our own standard for style, please defer to
the Python's PEP 7 \citep{pep7] for C style and emulating that in C++.

A basic set of guidelines to use when submitting code. Defer to the PEP7 standard with the following additions:

* All files and directories shall be lowercase and contain no spaces.
* Structs and class names should use `LeadingCapitals`.
* Enums and constants should be `BLOCK_CAPITALS`.
* Class members and functions without a parent class should be `snake_case`.
With the exception of when one of the words is a "class" name from
libembroidery in which case it has the middle capitals like this:
`embArray_add`.
* Don't use exceptions.
* Don't use ternary operator (?:) in place of if/else.
* Don't repeat a variable name that already occurs in an outer scope.

\subsection{Version Control

Being an open source project, developers can grab the latest code at any
time and attempt to build it themselves. We try our best to ensure that
it will build smoothly at any time, although occasionally we do break
the build. In these instances, please provide a patch, pull request
which fixes the issue or open an issue and notify us of the problem, as
we may not be aware of it and we can build fine.

Try to group commits based on what they are related to:
features/bugs/comments/graphics/commands/etc...

\subsection{Comments

When writing code, sometimes there are items that we know can be
improved, incomplete or need special clarification. In these cases, use
the types of comments shown below. They are pretty standard and are
highlighted by many editors to make reviewing code easier. We also use
shell scripts to parse the code to find all of these occurrences so
someone wanting to go on a bug hunt will be able to easily see which
areas of the code need more love.

libembroidery and Embroidermodder are written in C and adheres to C89 standards. This means
that any C99 or C++ comments will show up as errors when compiling with
gcc. In any C code, you must use:

----
/* Use C Style Comments within code blocks.
 *
 * Use Doxygen style code blocks to place todo, bug, hack, warning,
 * and note items like this:
 *
 * \todo EXAMPLE: This code clearly needs more work or further review.
 *
 * \bug This code is definitely wrong. It needs fixed.
 *
 * \hack This code shouldn't be written this way or I don't
 * feel right about it. There may a better solution
 *
 * \warning Think twice (or more times) before changing this code.
 * I put this here for a good reason.
 *
 * \note This comment is much more important than lesser comments.
 */
----

\section{Ideas

\subsection{Why this document

I've been trying to make this document indirectly through the Github
issues page and the website we're building but I think a
straightforward, plain-text file needs to be the ultimate backup for
this. Then I can have a printout while I'm working on the project.

\subsection{Qt and dependencies

I'm switching to SDL2 (which is a whole other conversation) which means
we can ship it with the source code package meaning only a basic build
environment is necessary to build it.

\subsection{Documentation

Can we treat the website being a duplicate of the docs a non-starter?
I'd be happier with tex/pdf only and (I know this is counter-intuitive)
one per project.

\subsection{Social Platform

So... all the issues and project boards etc. being on Github is all
well and good assuming that we have our own copies. But we don't if
Github goes down or some other major player takes over the space and we
have to move (again, since this started on SourceForge).

This file is a backup for that which is why I'm repeating myself between
them.

\subsection{Identify the meaning of these TODO items

* Saving CSV/SVG (rt) + CSV read/write UNKNOWN interpreted as COLOR bug `#179`
* Lego Mindstorms NXT/EV3 ports and/or commands

\subsection{Progress Chart

The chart of successful from-to conversions (previously a separate issue)
is something that should appear in the README.

\subsection{Standard

The criteria for a good Pull Request from an outside developer has these properties, from most to least important:

* No regressions on testing.
* Add a feature, bug fix or documentation that is already agreed on through
  GitHub issues or some other way with a core developer.
* No GUI specific code should be in libembroidery, that's for Embroidermodder.
* Pedantic/ansi C unless there's a good reason to use another language.
* Meet the style above (i.e.  PEP 7, Code Lay-out
  (https://peps.python.org/pep-0007/#code-lay-out}). We'll just fix the style
  if the code's good and it's not a lot of work.
* `embroider` should be in POSIX style as a command line program.
* No dependancies that aren't "standard", i.e. use only the C Standard Library.

\subsection{Image Fitting

A currently unsolved problem in development that warrants further research is
the scenario where a user wants to feed embroider an image that can then be .

\subsection{To Place

A _right-handed coordinate system_ index[right-handed coordinate system}
is one where up is positive and right is
positive. Left-handed is up is positive, left is positive. Screens often use
down is positive, right is positive, including the OpenGL standard so when
switching between graphics formats and stitch formats we need to use a vertical
flip (`embPattern_flip`).

`0x20` is the space symbol, so when padding either 0 or space is preferred and
in the case of space use the literal ' '.

\subsection{To Do

We currently need help with:

* Thorough descriptions of each embroidery format.
* Finding resources for each of the branded thread libraries (along with a full
  citation for documentation).
* Finding resources for each geometric algorithm used (along with a full
  citation for documentation).
* Completing the full `--full-test-suite` with no segfaults and at least a clear
  error message (for example ``not implemented yet'').
* Identifying ``best guesses'' for filling in missing information when going
  from, say `.csv` to a late `.pes` version. What should the default be when
  the data doesn't clarify?
* Improving the written documentation.
* Funding, see the Sponsor button above. We can treat this as ``work'' and put
  far more hours in with broad support in small donations from people who want
  specific features.

Beyond this the development targets are categories sorted into:

* Basic Features
* Code quality and user friendliness
* embroider CLI
* Documentation
* GUI
* electronics development

\subsection{Basic features}

* Incorporate `#if 0` ed parts of `libembroidery.c`.
* Interpret how to write formats that have a read mode from the source
  code and vice versa.
* Document the specifics of the file formats here for embroidery machine
  specific formats. Find websites and other sources that break down the binary
  formats we currently don't understand.
* Find more and better documentation of the structure of the headers for the
  formats we do understand.

\subsection{Code quality and user friendliness}

* Document all structs, macros and functions (will contribute directly
   on the web version).
* Incorporate experimental code, improve support for language bindings.
* Make stitch x, y into an EmbVector.

\subsection{Documentation}

Run `sloccount` on `extern/` and `.` (and ) so we know the
current scale of the project, aim to get this number low. Report the total as
part of the documentation.

Try to get as much of the source code that we maintain into C as possible so
new developers don't need to learn multiple languages to have an effect. This
bars the embedded parts of the code.

\subsection{GUI}

* Make EmbroideryMobile (Android) also backend to `libembroidery` with a Java wrapper.
* Make EmbroideryMobile (iOS) also backend to `libembroidery` with a Swift wrapper.
* Share some of the MobileViewer and iMobileViewer layout with the main EM2. Perhaps combine those 3 into the Embroidermodder repository so there are 4 repositories total.
* Convert layout data to JSON format and use cJSON for parsing.

\section{Development}

\subsection{Contributing}

If you're interested in getting involved, here's some guidance
for new developers. Currently The Embroidermodder Team is all
hobbyists with an interest in making embroidery machines more
open and user friendly. If you'd like to support us in some other way
you can donate to our Open Collective page (click the Donate button) so
we can spend more time working on the project.

All code written for libembroidery should be ANSI C89 compliant
if it is C. Using other languages should only be used where
necessary to support bindings.

\subsection{Debug}

If you wish to help with development, run this debug script and send us the
error log.

\begin{lstlisting}
#!/bin/bash

rm -fr libembroidery-debug

git clone http://github.com/embroidermodder/libembroidery libembroidery-debug
cd libembroidery-debug

cmake -DCMAKE_BUILD_TYPE=DEBUG .
cmake --build . --config=DEBUG

valgrind ./embroider --full-test-suite
\end{lstlisting}

While we will attempt to maintain good results from this script as part of
normal development it should be the first point of failure on any system we
haven't tested or format we understand less.

\subsection{Binary download}

We need a current `embroider` command line program download, so people can update
without building.

\chapter{Programming principles for the C core}

End arrays of char arrays with the symbol "END", the code will never require
this symbol as an entry.

Define an array as one of 3 kinds: constant, editable or data.

* Constant arrays are defined const and have fixed length matching the data.
* Editable arrays are fixed length, but to a length based on the practical use
  of that array. A dropdown menu can't contain more than 30 items, because we
  don't want to flood the user with options. However it can nest indefinately,
  so it is not restricted to a total number of entries.
* Data arrays is editable and changes total size at runtime to account for user data.

\section{Style rules for arrays}

1.


\chapter{Libembroidery on Embedded Systems}

The libembroidery library is designed to support embedded environments as well
as desktop, so it can be used in CNC applications.

Originally, the embedded system aspect of the Embroidermodder project was
targeted at the higher end \index{Arduino] prototyping board as part
of a general effort to make our own open source hardware (\index{OSHW]).

However, the task of building the interface for a full OSHW embroidery machine
neatly splits into the tasks of building a user interface to issue the
commands and the rig itself starting with the stepper motors that wire into
this control circuit. A well built control circuit could issue commands to
a variety of different machine layouts (for example many features are not
present on some machines)

\section{Compatible Boards}

We recommend using an \index{Arduino] Mega 2560 or another board with equal or
greater specs. That being said, we have had success using an Arduino Uno
R3 but this will likely require further optimization and other
improvements to ensure continued compatibility with the Uno. See below
for more information.

\section{Arduino Considerations}

There are two main concerns here: Flash Storage and SRAM.

libembroidery continually outgrows the 32KB of Flash storage on the
Arduino Uno and every time this occurs, a decision has to be made as to
what capabilities should be included or omitted. While reading files is
the main focus on arduino, writing files may also play a bigger role
in the future. Long term, it would be most practical to handle the
inclusion or omission of any feature via a single configuration header
file that the user can modify to suit their needs.

SRAM is in extremely limited supply and it will deplete quickly so any
dynamic allocation should occur early during the setup phase of the
sketch and sparingly or not at all later in the sketch. To help minimize
SRAM consumption on Arduino and ensure libembroidery can be used in any
way the sketch creator desires, it is required that any sketch using
libembroidery must implement event handlers. See the ino-event source
and header files for more information.

There is also an excellent article by Bill Earl on the Adafruit Learning
System
\footnote{\url{http://learn.adafruit.com/memories-of-an-arduino?view=all}}.

\section{Space}

Since a stitch takes 3 bytes of storage and many patterns use more than
10k stitches, we can't assume that the pattern will fit in memory. Therefore
we will need to buffer the current pattern on and off storage in small
chunks. By the same reasoning, we can't load all of one struct beore
looping so we will need functions similar to binaryReadInt16 for each
struct.

This means the EmbArray approach won't work since we need to load
each element and dynamic memory management is unnecessary because
the arrays lie in storage.

TODO: Replace EmbArray functions with `embPattern_` load functions.

\section{Tables}

All thread tables and large text blocks are too big to compile directly
into the source code. Instead we can package the library with a data packet
that is compiled from an assembly program in raw format so the specific
padding can be controlled.

In the user section above we will make it clear that this file
needs to be loaded on the pattern USB/SD card or the program won't function.

TODO: Start file with a list of offsets to data with a corresponding table
to load into with macro constants for each label needed.

\section{Current Pattern Memory Management}

It will be simpler to make one file per EmbArray so we keep an EmbFile*
and a length, so no malloc call is necessary. So there needs to be a consistent
tmpfile naming scheme.

TODO: For each pattern generate a random string of hexadecimal and append it
to the filenames like `stitchList_A16F.dat`. Need to check for a file
which indicates that this string has been used already.

\section{Special Notes}

Due to historical reasons and to remain compatible with the Arduino 1.0
IDE, this folder must be called ``utility''. Refer to the arduino build
process for more infofootnote:[https://arduino.github.io/arduino-cli/0.19/sketch-build-process/].

libembroidery relies on the Arduino SD library for reading files. See
the ino-file source and header files for more information.

\section{The Assembly Split}

One problem to the problem of supporting both systems with abundant memory
(such as a 2010s or later desktop) and with scarce memory (such as embedded
systems) is that they don't share the same assembly language. To deal with
this: there will be two equivalent software which are hand engineered to be
similar but one will be in C and the other in the assembly dialects we support.

All assembly will be intended for embedded systems only, since a slightly
smaller set of features will be supported. However, we will write a
`x86` version since that can be tested.

That way the work that has been done to simplify the C code can be applied
to the assembly versions.

\chapter{Electronics development}

\section{Ideas}

Currently experimenting with Fritzing 8 (8), upload netlists to embroiderbot when
they can run simulations using the asm in `libembroidery`.

Create a common assembly for data that is the same across chipsets
`libembrodiery_data_internal.s`.

Make the defines part of `embroidery.h` all systems and the function list
`c code only`. That way we can share some development between assembly and C versions.

\chapter{Mobile}

Again, it would help to use the C library we have already developed,
however for Android the supported platform is for Java applications.

\url{https://github.com/java-native-access/jna}

\url{https://github.com/marketplace/actions/setup-android-ndk}

\url{https://developer.android.com/ndk/guides}

See the bindings section for how this is achieved.

\bibliographystyle{usrtnat}
\bibliography{references}

\appendix

\chapter{Tread Tables}

\section{Arc Polyester Codes

.Arc Polyester Codes
[%header,format=csv}
\begin{tabular}{l l l}
include{data/shv_colors.csv}
\end{tabular}\section{Arc Rayon Codes

.Arc Rayon Codes
[%header,format=csv}
\begin{tabular}{l l l}
include{data/shv_colors.csv}
\end{tabular}
\section{Coats and Clark Rayon Codes

.Coats and Clark Rayon Codes
[%header,format=csv}
\begin{tabular}{l l l}
include{data/shv_colors.csv}
\end{tabular}
\section{DXF Colors

Based on the DraftSight color table.

.DXF Colors
[%header,format=csv}
\begin{tabular}{l l l}
include{data/shv_colors.csv}
\end{tabular}

\section{Exquisite Polyester Codes

.Exquisite Polyester Codes
[%header,format=csv}
\begin{tabular}{l l l}
include{data/shv_colors.csv}
\end{tabular}

\section{Fufu Polyester Codes

.Fufu Polyester Codes
[%header,format=csv}
\begin{tabular}{l l l}
include{data/shv_colors.csv}
\end{tabular}

\section{Fufu Rayon Codes

.Fufu Rayon Codes
[%header,format=csv}
\begin{tabular}{l l l}
include{data/shv_colors.csv}
\end{tabular}

\section{Hemingworth Polyester Codes

.Hemingworth Polyester Codes
[%header,format=csv}
\begin{tabular}{l l l}
include{data/shv_colors.csv}
\end{tabular}

\section{HUS Colors (find a citation)

; CSV format: red, green, blue, name, catalog_number

.HUS Colors
[%header,format=csv}
\begin{tabular}{l l l}
include{data/shv_colors.csv}
\end{tabular}

\section{Isacord Polyester Codes

.Isacord Polyester Codes
[%header,format=csv}
\begin{tabular}{l l l}
include{data/shv_colors.csv}
\end{tabular}

\section{Isafil Rayon Codes

.Isafil Rayon Codes
[%header,format=csv}
\begin{tabular}{l l l}
include{data/shv_colors.csv}
\end{tabular}

\section{JEF Colors

To do: find a citation

.JEF Colors
[%header,format=csv}
\begin{tabular}{l l l}
include{data/shv_colors.csv}
\end{tabular}

\section{Madeira Polyester Codes

.Madeira Polyester Codes
[%header,format=csv}
\begin{tabular}{l l l}
include{data/shv_colors.csv}
\end{tabular}

\section{Madeira Rayon Codes

.Madeira Rayon Codes
[%header,format=csv}
\begin{tabular}{l l l}
include{data/shv_colors.csv}
\end{tabular}

\section{Marathon Polyester Codes

.Marathon Polyester Codes
[%header,format=csv}
\begin{tabular}{l l l}
include{data/shv_colors.csv}
\end{tabular}

\section{Marathon Rayon Codes

.Marathon Rayon Codes
[%header,format=csv}
\begin{tabular}{l l l}
include{data/shv_colors.csv}
\end{tabular}

\section{Metro Polyester Code

.Metro Polyester Code
[%header,format=csv}
\begin{tabular}{l l l}
include{data/shv_colors.csv}
\end{tabular}

\section{Pantone Codes

See https://www.pantone-colours.com/}

.Pantone Codes
[%header,format=csv}
\begin{tabular}{l l l}
include{data/shv_colors.csv}
\end{tabular}

\section{PCM Color Codes

.PCM Color Codes
[%header,format=csv}
\begin{tabular}{l l l}
include{data/shv_colors.csv}
\end{tabular}

\section{PEC Color Codes

.PEC Color Codes
[%header,format=csv}
\begin{tabular}{l l l}
include{data/shv_colors.csv}
\end{tabular}

\section{Robinson Anton Polyester Codes
index[Robinson Anton}

.Robinson Anton Polyester Colors.
[%header,format=csv}
\begin{tabular}{l l l}
include{data/shv_colors.csv}
\end{tabular}

\begin{table}
\begin{tabular}{l l l}
\include{data/shv_colors.csv}
\end{tabular}
\caption{
Robinson Anton Rayon Colors.
\index{Robinson Anton}
\index{Rayon}
}
\end{table}

\begin{table}
\begin{tabular}{l l l}
\include{data/shv_colors.csv}
\end{tabular}
\caption{SHV Thread Colors.}
\end{table}

\section{Sigma Polyester Codes}

.Sigma Polyester Thread Colors
[%header,format=csv}
\begin{tabular}{l l l}
\include{data/z102_isacord_polyester_colors.csv}
\end{tabular}

\section{Sulky Rayon Colors
\index{Sulky}
\index{Rayon}

.Sulky Rayon Colors
[%header,format=csv}
\begin{tabular}{l l l}
include{data/z102_isacord_polyester_colors.csv}
\end{tabular}

\section{SVG Colors}
\index{SVG}

Converted from the table at:
\url{https://www.w3.org/TR/SVGb/types.html#ColorKeywords}

NOTE: This supports both UK and US English names, so the repeated values aren't an error.

[%header,format=csv}
\begin{tabular}{l l l}
include{data/z102_isacord_polyester_colors.csv}
\end{tabular}

\section{ThreadArt Polyester
index[ThreadArt}

\begin{table}
\begin{tabular}{l l l}
include{data/z102_isacord_polyester_colors.csv}
\end{tabular}
\caption{.}
\end{table}

\section{ThreadArt Rayon Codes}
\index[ThreadArt}

\begin{table}
\begin{tabular}{l l l}
\include{data/z102_isacord_polyester_colors.csv}
\end{tabular}
\caption{.}
\end{table}

\section{ThreaDelight Polyester Codes}
\index[ThreadDelight}

\begin{table}
\begin{tabular}{l l l}
\include{data/z102_isacord_polyester_colors.csv}
\end{tabular}
\caption{.}
\end{table}

\section{Z102 Isacord Polyester Codes}

\begin{table}
\begin{tabular}{l l l}
\include{data/z102_isacord_polyester_colors.csv}
\end{tabular}
\caption{.}
\end{table}

\printindex
